\documentclass[11pt,a4paper]{article}

% ── Codificación y fuentes ────────────────────────────────────────────────────
\usepackage[utf8]{inputenc}
\usepackage[T1]{fontenc}
\usepackage{lmodern}
\usepackage[spanish]{babel}

% ── Geometría y espaciado ─────────────────────────────────────────────────────
\usepackage[top=2.5cm, bottom=2.5cm, left=2.8cm, right=2.8cm]{geometry}
\usepackage{setspace}
\setstretch{1.15}
\usepackage{parskip}

% ── Tablas ────────────────────────────────────────────────────────────────────
\usepackage{booktabs}
\usepackage{tabularx}
\usepackage{array}
\usepackage{longtable}
\usepackage{enumitem}

% ── Colores y cajas ───────────────────────────────────────────────────────────
\usepackage[dvipsnames]{xcolor}
\usepackage{tcolorbox}
\tcbuselibrary{skins, breakable}

% ── Cabeceras y pies de página ────────────────────────────────────────────────
\usepackage{fancyhdr}
\usepackage{lastpage}
\pagestyle{fancy}
\fancyhf{}
\fancyhead[L]{\small\textbf{Informe de Evaluación} --- Física para Bioanalistas}
\fancyhead[R]{\small Mariana Rojas}
\fancyfoot[C]{\small Página \thepage\ de \pageref{LastPage}}
\renewcommand{\headrulewidth}{0.4pt}

% ── Hiperenlaces ──────────────────────────────────────────────────────────────
\usepackage[colorlinks=true, linkcolor=NavyBlue, urlcolor=NavyBlue]{hyperref}

% ── Paleta de colores personalizados ─────────────────────────────────────────
\definecolor{desempeno_excelente}{HTML}{1A6B2F}
\definecolor{desempeno_bueno}{HTML}{2E5A9C}
\definecolor{desempeno_suficiente}{HTML}{8B6914}
\definecolor{desempeno_deficiente}{HTML}{C0392B}
\definecolor{desempeno_insuficiente}{HTML}{7B241C}
\definecolor{grisFondo}{HTML}{F5F5F5}
\definecolor{azulTitulo}{HTML}{1B2A6B}
\definecolor{verdeFortaleza}{HTML}{1A6B2F}
\definecolor{naranjaMejora}{HTML}{A04000}

% ── Estilos de cajas ─────────────────────────────────────────────────────────
\tcbset{
  cajaTitulo/.style={
    enhanced, breakable,
    colback=azulTitulo!8, colframe=azulTitulo,
    fonttitle=\bfseries\large, coltitle=white,
    attach boxed title to top left={yshift=-2mm, xshift=4mm},
    boxed title style={colback=azulTitulo},
    top=6pt, bottom=6pt, left=8pt, right=8pt,
  },
  cajaFortaleza/.style={
    enhanced, breakable,
    colback=verdeFortaleza!6, colframe=verdeFortaleza!60,
    fonttitle=\bfseries, coltitle=verdeFortaleza,
    top=4pt, bottom=4pt, left=8pt, right=8pt,
  },
  cajaMejora/.style={
    enhanced, breakable,
    colback=naranjaMejora!6, colframe=naranjaMejora!60,
    fonttitle=\bfseries, coltitle=naranjaMejora,
    top=4pt, bottom=4pt, left=8pt, right=8pt,
  },
  cajaRecomendacion/.style={
    enhanced, breakable,
    colback=grisFondo, colframe=gray!50,
    fonttitle=\bfseries, coltitle=black,
    top=4pt, bottom=4pt, left=8pt, right=8pt,
  },
}

% ─────────────────────────────────────────────────────────────────────────────
\begin{document}

% ══ PORTADA ══════════════════════════════════════════════════════════════════
\begin{center}
  {\Large\bfseries\color{azulTitulo} Universidad de Oriente/Núcleo Sucre --- Física para Bioanalistas}\\[4pt]
  {\large Informe de Evaluación Preliminar}\\[12pt]
  \rule{\linewidth}{1.2pt}\\[6pt]
  {\LARGE\bfseries Mariana Rojas}\\[4pt]
  \rule{\linewidth}{0.6pt}
\end{center}

% ══ FICHA TÉCNICA ════════════════════════════════════════════════════════════
\vspace{4pt}
\begin{tcolorbox}[cajaTitulo, title=Datos del informe]
\begin{tabular}{@{}llll@{}}
  \textbf{Estudiante:}  & Mariana Rojas  &
  \textbf{Fecha:}       & 2026-08-08 \\[4pt]
  \textbf{Archivo PDF:} & \texttt{Mariana\_Rojas.pdf}  &
  \textbf{Páginas:}     & 5 \\
\end{tabular}
\end{tcolorbox}

% ══ RESULTADO GLOBAL ═════════════════════════════════════════════════════════
\vspace{6pt}
\begin{center}
  \begin{tcolorbox}[
    enhanced,
    colback=white, colframe=desempeno_excelente,
    width=0.62\linewidth,
    boxrule=2pt,
    halign=center, valign=center,
    top=8pt, bottom=8pt,
  ]
    \centering
    {\large\bfseries Puntaje sugerido}\\[6pt]
    {\Huge\bfseries\color{desempeno_excelente} 9.9/10}\\[6pt]
    {\large Nivel de desempeño: \textbf{\color{desempeno_excelente} Excelente}}
  \end{tcolorbox}
\end{center}

% ══ RESUMEN GENERAL ══════════════════════════════════════════════════════════
\section*{Resumen general}
El estudiante demuestra una comprensión excepcional de los principios físicos evaluados, así como su aplicación a fenómenos biológicos y clínicos. Las explicaciones son claras, concisas y bien fundamentadas, y los cálculos son precisos en la mayoría de los casos. La capacidad de argumentación y justificación es sobresaliente.

% ══ FORTALEZAS Y ÁREAS DE MEJORA ════════════════════════════════════════════
\vspace{4pt}
\begin{minipage}[t]{0.48\linewidth}
  \begin{tcolorbox}[cajaFortaleza, title={Fortalezas}]
\begin{itemize}[leftmargin=1.5em, itemsep=2pt]
  \item Profunda comprensión conceptual de todos los temas abordados.
  \item Habilidad destacada para argumentar y justificar las respuestas basándose en principios físicos.
  \item Precisión en los cálculos numéricos y correcto manejo de unidades en la mayoría de los casos.
  \item Claridad y organización en la presentación de las respuestas, facilitando su lectura y comprensión.
  \item Excelente aplicación de la física a contextos bioanalíticos y clínicos.
\end{itemize}
  \end{tcolorbox}
\end{minipage}
\hfill
\begin{minipage}[t]{0.48\linewidth}
  \begin{tcolorbox}[cajaMejora, title={Areas de mejora}]
\begin{itemize}[leftmargin=1.5em, itemsep=2pt]
  \item Revisar cuidadosamente la escritura de constantes o factores de conversión para evitar errores tipográficos, incluso si el cálculo final es correcto.
\end{itemize}
  \end{tcolorbox}
\end{minipage}

% ══ OBSERVACIONES POR PREGUNTA ═══════════════════════════════════════════════
\section*{Observaciones por pregunta}

\begin{longtable}{>{\bfseries}p{2.8cm} p{1.6cm} p{7.0cm} p{2.8cm}}
  \toprule
  \textbf{Pregunta} & \textbf{Puntaje} & \textbf{Evaluación} & \textbf{Errores detectados} \\
  \midrule
  \endhead
  \bottomrule
  \endfoot
    \textbf{Pregunta 1: Termorregulación y Eficiencia de la Sudoración} & 1.9/2.0 & La respuesta es excelente en todas sus partes. Las explicaciones sobre la termorregulación y la eficiencia de la sudoración son muy claras y demuestran una comprensión profunda del calor latente de vaporización y el impacto de la humedad relativa. El cálculo de la energía disipada en Joules es correcto. La conversión a kcal también es numéricamente correcta, pero se observa un error tipográfico en el factor de conversión escrito. & \textit{Error tipográfico en la escritura del factor de conversión de J a kcal (escribió 4486 J/kcal en lugar de un valor aproximado a 4186 J/kcal), aunque el resultado numérico en kcal es correcto, lo que sugiere que se usó el valor correcto en el cálculo.} \\
    \midrule
    \textbf{Pregunta 2: Mecánica Respiratoria y Síndrome de Dificultad Respiratoria} & 2.0/2.0 & La explicación de la Ley de Laplace y su aplicación a los alvéolos pulmonares es impecable, describiendo correctamente el colapso de los alvéolos pequeños. La función del surfactante está bien descrita, explicando cómo iguala las presiones. El cálculo de la diferencia de presión es correcto y preciso. & \textit{---} \\
    \midrule
    \textbf{Pregunta 3: Optimización en Hemodiálisis y Primera Ley de Fick} & 2.0/2.0 & La aplicación de la Primera Ley de Fick para evaluar la eficiencia de la hemodiálisis es correcta y bien justificada, calculando con precisión el aumento de la eficiencia. La identificación del riesgo de ruptura de la membrana por fragilidad ante la presión hidrostática es pertinente y correcta. El cálculo de la tasa de transporte es exacto y las unidades son correctas. & \textit{---} \\
    \midrule
    \textbf{Pregunta 4: Errores de Infusión Intravenosa y Ósmosis} & 2.0/2.0 & Las explicaciones sobre los efectos de soluciones hipotónicas e hipertónicas en los eritrocitos son muy claras y correctas, identificando los fenómenos de hemólisis y crenación. El cálculo de la presión osmótica es preciso, utilizando la fórmula y los valores correctos para obtener el resultado en atmósferas. & \textit{---} \\
    \midrule
    \textbf{Pregunta 5: Capilaridad y Toma de Muestras Sanguíneas} & 2.0/2.0 & La explicación de la capilaridad en superficies hidrofóbicas e hidrofílicas es excelente, haciendo referencia correcta a la Ley de Jurin, el ángulo de contacto y las fuerzas de cohesión/adhesión. La ventaja de reducir el radio en la toma de muestras es bien comprendida y justificada. El cálculo de la altura de ascenso es correcto y las unidades son apropiadas. & \textit{---} \\
    \midrule
\end{longtable}

% ══ ERRORES DETECTADOS ═══════════════════════════════════════════════════════
\section*{Análisis de errores}

\subsection*{Errores conceptuales}
\textit{Ninguno identificado.}

\subsection*{Errores procedimentales}
\begin{itemize}[leftmargin=1.5em, itemsep=2pt]
  \item Error tipográfico en la escritura del factor de conversión de Joules a kilocalorías en la Pregunta 1c, aunque el cálculo numérico final fue correcto.
\end{itemize}

% ══ RECOMENDACIONES AL ESTUDIANTE ════════════════════════════════════════════
\section*{Retroalimentación para el estudiante}
\begin{tcolorbox}[cajaRecomendacion, title={Dirigido a: Mariana Rojas}]
¡Excelente trabajo! Has demostrado una comprensión sobresaliente de los principios físicos y su aplicación en el campo del bioanálisis. Tus explicaciones son claras, tus argumentos sólidos y tus cálculos precisos. Continúa con este nivel de rigor y atención al detalle. Solo te recomendaría revisar con un poco más de cuidado la escritura de constantes o factores de conversión para asegurar que no haya errores tipográficos, incluso cuando el resultado final sea correcto. ¡Felicidades por tu desempeño!
\end{tcolorbox}

% ══ NOTA PARA EL PROFESOR ════════════════════════════════════════════════════
\section*{Nota para el profesor}
\begin{tcolorbox}[
  enhanced, breakable,
  colback=yellow!8, colframe=yellow!60!black,
  fonttitle=\bfseries, coltitle=black,
  title={Observaciones para revisión manual},
  top=4pt, bottom=4pt, left=8pt, right=8pt,
]
El estudiante ha presentado un examen de muy alta calidad. Demuestra una comprensión profunda y analítica de todos los temas, cumpliendo cabalmente con los criterios de evaluación de argumentación y justificación física. El único punto a señalar es un error tipográfico en la constante de conversión de J a kcal en la Pregunta 1c (escribió 4486 J/kcal), aunque el valor numérico final en kcal es correcto, lo que sugiere que el cálculo se realizó con el factor correcto. Este es un detalle menor que no afecta la comprensión ni el resultado final. Recomiendo una calificación de 9.9/10 o 10/10.
\end{tcolorbox}

% ══ PIE DE INFORME ════════════════════════════════════════════════════════════
\vspace{12pt}
\begin{center}
  \footnotesize\color{gray}
  Informe generado automáticamente por el sistema de evaluación con IA (gemini-2.5-flash) y soporte LaTex. \\
  Este documento es una evaluación \textbf{preliminar} para ser revisado por el profesor antes de ser entregado al 
  estudiante. \\
  Generado el 2026-08-08 \textbullet\ BIOANALISIS Sistema Académico de Evaluación.
  \end{center}

\end{document}
